\subsection{Wallet Integration}

The proposed system adopts a non-custodial wallet model in which students retain full control of their blockchain identities. Instead of generating wallets on behalf of users, the platform integrates with MetaMask, a browser-based Ethereum wallet extension.\\

When a student accesses the portal, the application requests authorization through the browser's injected Ethereum provider. Upon approval, the selected wallet address is retrieved and displayed within the user interface. The captured address is then submitted to the backend and associated with the student's academic profile. The wallet registration workflow is shown below:

\begin{center}
    \begin{tabular}{c}
        Student Portal           \\
        $\downarrow$             \\
        MetaMask Extension       \\
        $\downarrow$             \\
        Wallet Address Retrieval \\
        $\downarrow$             \\
        Backend API              \\
        $\downarrow$             \\
        Student Database Profile
    \end{tabular}
\end{center}

Only public wallet addresses are stored within the database. Private keys never leave the student's MetaMask wallet and remain under the student's exclusive control. After wallet registration, the stored address becomes the destination for NFT certificate issuance. During the graduation process, the backend retrieves the student's registered wallet address and invokes the smart contract minting function through the university's administrative blockchain account. The newly created NFT certificate is therefore transferred directly to the student's wallet. This architecture enables decentralized certificate ownership while eliminating the need for the institution to manage or store user private keys.